\section{System and Demonstration}\label{sec:system}

\textbf{Single entry point.} The command line, the batch mode and the interface all call one function,
\texttt{detect()}. It loads the image, applies the protocol, runs the eight modules with the tuned parameters
(Section~\ref{sec:params}), classifies with the calibrated forest, estimates the forgery type, localises with the
learned fusion and gates the mask by the verdict. It refuses a model whose recorded feature parameters or protocol
differ from the current configuration, so no entry point can drift from the evaluated system.

\textbf{Explanations.} Each result carries the exact SHAP contributions (TreeExplainer)~\cite{lundberg2017shap,
lundberg2020tree} of the 82 features to the forest's score, translated into plain language; a unit test checks that
they plus the base value reproduce the forest's probability exactly. The calibrated probability is a monotone
transformation of that score, so the directions carry over but the sizes are not on the probability scale. The verdict
text states the uncertain band (P(tampered) within $0.50\pm0.24$ under R gives no verdict) and the share of test
images in it (48\% under R, 4\% under B); evidence sentences name the strongest evidence and the known false-alarm
cause, repeated real structure mimicking a copy-move (Section~\ref{sec:test}). A PDF report and a JSON record can be
exported.

\textbf{Input preparation.} Every image is validated and normalised before analysis (Appendix~\ref{app:system}).

\begin{figure}[!t]
\centering
\includegraphics[width=\columnwidth, trim=375 1235 75 340, clip]{fig_ui_analyze.png}
\caption{Result panel of the ``Analyze an image'' page for a MICC-F220 copy-move example under protocol R: verdict
with P(tampered), likely type and size of the suspected region, next to the analysed image, the suspected region (red)
and the fused tampering probability map. Evidence sentences and the per-image SHAP chart follow below it on the page (not shown).}
\label{fig:ui-analyze}
\end{figure}

\textbf{Interface.} A four-page Streamlit application offers (1)~\emph{Analyze an image} (Fig.~\ref{fig:ui-analyze}):
verdict, likely type, suspected region, fused map, explanation, per-technique evidence maps and downloads, with the
protocol chosen in the sidebar (R by default); (2)~\emph{batch analysis} of up to 200 files or a folder, where an
unreadable file does not stop the batch; (3)~a \emph{dataset explorer} whose leak tab shows format and quantisation
tables by class and the shortcut classifiers' balanced accuracy per protocol; (4)~a \emph{results dashboard} that
reads the saved evaluation without recomputing it.

\textbf{Keeping the evaluation honest.} The explorer never lists the test split. CASIA examples are labelled as
training data, since the deployed models were refit on them; MICC-F220 images serve as unseen examples. The home page
states the limits: trained on CASIA, a verdict is evidence, not proof, about half of all images fall in the uncertain
band, and the known weaknesses.

\textbf{Performance and verification.} A CASIA-sized image is analysed in about 0.2~s, and every page is covered by
automated headless tests (Appendix~\ref{app:system}).
